\section{Agent Paradigm}
\label{sec:agent_paradigm}

The \gls{rl} framework is built around the agent paradigm, in which the decision-making entity is referred to as the agent.
A standard definition is:
\begin{quote}
\emph{``An agent is anything that can be viewed as perceiving its environment through sensors and acting upon that environment through actuators.''} \cite{russell2021artificial}
\end{quote}

When this definition is adapted to the \gls{rl} setting, it can be described as follows.
The agent perceives its environment through observations.
These observations are then processed by a decision-making mechanism referred to as the policy.
The policy determines how the agent modifies its environment through the actions it selects.
The resulting action modifies the environment and, in turn, affects its subsequent dynamics.

In the context of \gls{rl}, the policy is obtained through learning.
The objective of the agent is to maximize the cumulative reward it receives over time through its interactions with the environment.
The procedure responsible for this process is the \gls{rl} algorithm. The algorithm progressively modifies the policy based on the experience collected through interaction.

In this manuscript, the \gls{rl} algorithm is considered as a process external to the agent.
This distinction is a deliberate conceptual choice.
Some other presentations of \gls{rl}, describe the agent as both the learner and the decision-making entity \cite{suttonbarto1998}.
We chose to distinguish the agent from the \gls{rl} algorithm because this provides a more general definition of an agent.

An agent can therefore exist without learning, for example when its behavior is entirely specified by a set of programmed rules.
Such an agent still perceives, processes information, and acts upon its environment, but its behavior is not the result of \gls {rl}.